% La mirada geométrica
\section{La mirada geométrica}\label{cap:03}

El hilo conductor de la materia no es memorizar matrices, sino aprender a \emph{ver} qué
transformación describe cada fórmula. Antes de calcular conviene interrogar la figura.

\begin{center}
\begin{tabularx}{\textwidth}{>{\bfseries}p{38mm} X}
\toprule
Pregunta geométrica & Qué revela\\
\midrule
¿Qué queda fijo? & centro, eje, recta o dirección privilegiada\\
¿Cómo se mueve un punto? & trayectoria recta, circular o curva\\
¿Qué se conserva? & distancia, ángulo, paralelismo, razón, área u orientación\\
¿Qué cambia? & escala, dirección, forma, curvatura o sistema de referencia\\
¿Qué pasa con una figura simple? & la imagen de un segmento, triángulo o paralelogramo
determina la transformación\\
¿Qué hace la composición? & permite construir transformaciones complejas a partir de
operaciones simples\\
\bottomrule
\end{tabularx}
\end{center}

\begin{idea}[Regla de lectura]
Para cada teorema se seguirá este orden:
\begin{enumerate}
  \item \textbf{Construcción}: dibujar el punto y su imagen.
  \item \textbf{Movimiento}: describir la trayectoria y el sentido.
  \item \textbf{Invariantes}: identificar qué no cambia.
  \item \textbf{Composición}: seguir los pasos intermedios en el orden correcto.
  \item \textbf{Traducción algebraica}: escribir la matriz o fórmula que resume el dibujo.
  \item \textbf{Verificación}: usar un ejemplo numérico para comprobar, no para sustituir,
  la comprensión geométrica.
\end{enumerate}
\end{idea}

\begin{resultado}[La matriz como geometría comprimida]
Las columnas de una matriz \(A\) son las imágenes de los vectores de la base:
\[
  A=
  \begin{pmatrix}\vert&\vert\\Ae_1&Ae_2\\\vert&\vert\end{pmatrix}.
\]
Por eso, para interpretar geométricamente una transformación lineal basta observar qué hace
con dos direcciones independientes y con el paralelogramo que generan.
\end{resultado}


\subsection{Los tres objetos que no hay que confundir: puntos, reglas de medir y grillas}\label{sec:puntos-formas-grillas}

Síntesis mía, como brújula para el resto del apunte: en la materia circulan tres objetos
que se escriben con números parecidos y son cosas distintas. El hilo conductor es que
todo vector es un mojón o un desplazamiento entre mojones, y toda matriz es una manera de
mover el espacio o de cambiar la grilla con que se lo mide.

Los dos ladrillos de abajo de esa tabla, desarrollados en la subsección~\ref{sec:desaprender}
y en la Unidad 0, son el idioma en que están escritas las tres filas:

\begin{description}
  \item[Grupos de transformaciones (ladrillo 1).] Las transformaciones de un mismo tipo
  (rotaciones, traslaciones, afinidades, homografías) forman un grupo: componer dos da otra
  del mismo tipo, la identidad es el neutro y cada una tiene inversa. ``Transformaciones y
  Operaciones [\ldots] son dos conceptos que van atravesando la materia'' (resumen 2025,
  clase 1); la inversa ``es un elemento de un conjunto que, al operar con otro resulta en el
  elemento neutro respecto de esa operación'' (resumen 2025, clase 10). Cuándo conmuta el
  grupo es la pregunta que decide si se puede armar un sistema de coordenadas: las
  rotaciones del plano y las traslaciones conmutan; las rotaciones de \(\R^3\) no
  (sección~\ref{sec:orden-importa}). En NumPy la operación es \texttt{A @ B}, el neutro
  \texttt{np.eye(n)} y la inversa aplicada a un dato \texttt{np.linalg.solve(A, y)}.
  \item[Los cuerpos \(\R\) y \(\mathbb{C}\) como matrices (ladrillo 2).] Los escalares
  escalan: un real \(a\) actúa sobre la recta como \(x\mapsto ax\), con inversa \(1/a\) y
  determinante \(|a|\). Un complejo \(z=a+ib\) actúa sobre el plano como la matriz
  \[
    M(z)=\begin{pmatrix}a&-b\\ b&a\end{pmatrix},\qquad
    M(z_1+z_2)=M(z_1)+M(z_2),\quad M(z_1z_2)=M(z_1)M(z_2),\quad M(1/z)=M(z)^{-1},
  \]
  es decir rota \(\arg z\) y escala \(|z|\) a la vez, con \(\det M(z)=|z|^2\): las matrices
  de esa forma \emph{son} \(\mathbb{C}\) con otro ropaje. ``La matriz de rotación de 2x2 que
  vimos en algebra del CBC hace exactamente la misma operación que hacíamos con los
  complejos'' (resumen 2025, clase 1; la equivalencia de estructuras la remite a Gentile,
  \emph{Complejos y Generalización}, apéndice~\ref{ap:bibliografia}). Sección~\ref{sec:rot-complejos}.
\end{description}

Con el grupo y el cuerpo puestos, los tres objetos:

{\small
\begin{tabularx}{\textwidth}{>{\bfseries\raggedright\arraybackslash}p{27mm}>{\raggedright\arraybackslash}p{30mm}YY}
\toprule
Objeto & Cómo se escribe, dónde vive & Qué es geométricamente & Cómo se transforma\\
\midrule
Punto / vector & columna \(v\); vive en el espacio \(V\) &
una flecha: un mojón (desde el origen) o un desplazamiento entre dos mojones &
con la matriz, por la izquierda: \(v\mapsto Av\)\\
Forma lineal / gradiente & fila \(\omega\); vive en el dual \(V^{*}\) &
un mazo de rectas (en \(\R^3\), planos) de nivel equiespaciadas: una regla graduada
tirada en el plano; evalúa un vector \emph{sin} producto interno, mirando qué recta del
mazo pasa por él &
del lado contrario: \(\omega\mapsto\omega A^{-1}\), para que \(\omega v\) no cambie\\
Producto interno / métrica & matriz simétrica \(G\) (en una superficie, la primera forma
fundamental \(E,F,G\)) &
la grilla con que se mide: permite largos y ángulos, \(\|v\|^2=v^TGv\); es el único
puente entre una forma lineal y una dirección perpendicular (``el gradiente es
perpendicular a la curva de nivel'' necesita a \(G\)) &
como forma cuadrática: \(G\mapsto J^TGJ\) (sección~\ref{sec:metrica})\\
\bottomrule
\end{tabularx}
}

Sin producto interno no hay ``vector perpendicular a una recta'': hay la recta (forma) y
hay los puntos (vectores), y nada más. Con producto interno, la forma \(\omega\) se
``disfraza'' del vector \(G^{-1}\omega^T\) y recién ahí se puede hablar de la dirección
normal. En la base canónica \(G=I\) y el disfraz es transponer la fila, por eso en el
plano la confusión no cuesta nada; en una superficie, donde \(E,F,G\) cambian de punto a
punto, cuesta metros.

\begin{atencion}[Por qué el gradiente no es un vector]
\(df=\nabla f\,dx\) es una forma lineal: toma un desplazamiento \(dx\) y devuelve cuánto
cambia \(f\); sus curvas de nivel son el plano tangente a la gráfica, y su recta de nivel
cero es la tangente a la curva de nivel de \(f\). Prueba de que no es un vector: si
estiro el eje \(x\) al doble, los puntos duplican su primera coordenada, pero las marcas
de la regla \(\nabla f\) a lo largo de \(x\) se separan a la mitad (hay que avanzar el
doble para que \(f\) suba lo mismo), es decir la componente se divide por dos: se
transforma al revés que un vector. El profesor se corrigió en vivo el 14 de octubre de
2026, mientras resolvía el plano normal a una curva intersección de dos superficies: ``el
gradiente de una función es un vector\ldots{} no, mentí. Lo dije mal. [\ldots] El gradiente
de la función es una forma lineal, no un vector, porque me permite hacer una aproximación
lineal de la función'' (transcripción auto-traducida, \file{fuentes/transcripciones/2026-10-14.md}).
Y es lo que el resumen 2025 de la clase 9 llama ``dato clave: el gradiente es una forma
lineal no un vector''. Lo que sí es un vector es la normal, que se obtiene con el producto
vectorial de dos tangentes, o levantando el gradiente con la métrica.
\end{atencion}

\begin{resultado}[El puente al cálculo en superficies]
Los cinco ladrillos de la subsección~\ref{sec:desaprender} son también el idioma de las
clases 9 a 14. Derivar una parametrización \(r(u,v)\) es usar una matriz lineal, la
jacobiana \(J=(r_u\ \ r_v)\), en el espacio tangente: entra un pasito \((du,dv)\) de la
grilla, sale el pasito \(J\,(du,dv)^T\) sobre la superficie (ladrillos 1 y 3,
sección~\ref{sec:variedades}). El gradiente de una función escalar es una forma del dual,
\(df=\nabla f\,dx\) (ladrillo 4). La primera forma fundamental \(J^TJ=
\bigl(\begin{smallmatrix}E&F\\ F&G\end{smallmatrix}\bigr)\) es la métrica del tangente,
la única cinta con que se miden largos y ángulos en la grilla \((u,v)\) (ladrillo 5,
sección~\ref{sec:metrica}); por eso la perpendicular dibujada en \((\varphi,\lambda)\) no
es la del terreno, y el gradiente levantado con esa métrica, \(G^{-1}\nabla f^{T}\), sí lo
es. Y curvarse, la hessiana o la segunda forma fundamental, es una forma cuadrática: un
número por cada dirección, cuya diagonalización da las direcciones principales
(sección~\ref{sec:curvatura}).
\end{resultado}
